\scp{\tsc{East Sumbawa}}{03-01}
The exact number of languages/dialects spoken in the eastern
part of Sumbawa Island (corresponding to the current Indonesian
regencies of Bima and Dompu) remains to be properly investigated.
At least three languages are identifiable: Bima, Kolo, and Sanggar.
The largest and most well described is Bima [bhp] (a.k.a.
Mbojo).\footnote{
		In May 2015, four Mbojo speakers (both male
		and female) from East Sumbawa met briefly with Grimes
		and explained that their language was different from Bima spoken in the city.
		It may be both a different dialect
		and an alternate name for the language.
		The exact basis of the difference is unclear,
		but the idea of “Mbojo” being an alternate name,
		and also a dialect/variety of Bima needs to be pursued further.}
Sanggar is spoken on the eastern part of the Sanggar peninsula.
At the time of writing,
it has not been assigned an ISO 693-3 code,
but the lexical data in \citet{Zollinger1850} shows that it is
clearly a separate language from Bima.\footnote{
		Thanks to Harald Hammarström for bringing the existence of Sanggar
		to our attention.} 

In addition to these two languages,
\citet[3--4]{Jonker1896} reports the existence of two “dialects” of
Bima: Kolo (spoken in the eponymous village on the north-eastern
part of the bay of Bima \citep{PartamiKaryawan1995}, and Toloweri.
Although available Kolo data is limited,
it has differences in sound changes,
functors, basic vocabulary, and morphology compared with Bima.
Thus, we separate it off as a separate language.\footnote{
		\citet[4]{Jonker1896}
		writes: “Kolo can be considered a dialect of Bima.” However,
		Jonker uses the term “dialect” for speech varieties which are
		separate languages, though more closely related to one another
		than to other neighbouring languages.
		Thus, for instance, all the different languages of Rote Island
		(\srf{04-02}) are labelled “dialects” by Jonker.}
Jonker reports
that “The language of Tolo-Weri corresponds completely to that
of Kolo with regard to sound changes,
and to a great extent also with regard to vocabulary”.\footnote{
“De
taal van Tòlo-Weri komt volkomen overeen met die van Kòlo ten
opzichte der klankverschijnselen en grootendeels ook wat betreft
den woordenschat,” \citep[4]{Jonker1896}.}
However,
almost no Toloweri data is available to make an independent
assessment of its relationship with Kolo or Bima.

There are a number of changes that are shared between Bima,
Kolo, and Sanggar which allow them to be placed within an \tsc{East
Sumbawa} group: *R > {\0} in most words, *-ay > \it{e,}
and *-aw > \it{o}.
All of these sound changes also occur in Komodo (\srf{03-01-01}),
as well as certain other \tsc{Bima-Lembata} languages,
thus considerably weakening the evidence for \tsc{East Sumbawa}
as a distinct subgroup within \tsc{Bima-Lembata}.

Among the \tsc{East Sumbawa} languages,
Bima shows a large number of unconditioned splits,
most of which are summarised by \citet[75--83]{Blust2008}.
The following unconditioned splits occur in Bima: *p > \it{f,
p}; *t > \it{d,
t}; *k > \it{h,
k}; initial *b- > \it{ɓ,
b, mb, w}; initial *d- > \it{r,
ɗ}; *j/*l > \it{r,
l}; *z > \it{r,
ɗ}; *s > \it{s,
ʧ}; *R > \it{r,} {\0};
and *w > \it{w,} {\0}.
Some of these splits have partial conditioning environments,
but exceptions to these environments also occur.
For the most part these splits are not shared with Sanggar or Kolo.
The changes *t > \it{d,
t}, *l > \it{r, l},
and (perhaps) *p > \it{f,
p} and *s > \it{s,
ʧ} are shared with \tsc{Havu-Dhao} (\srf{03-03-03}).
Similarly, there are some cases of medial *k > \it{h} shared
between Bima and Komodo.

There is also disagreement between \citet{Zollinger1850}
and more modern sources as to whether a certain word attests
a particular sound change in Bima.
Four examples which appear unlikely to be due to error are the
following (words from \citet{Zollinger1850} are given first): *saŋan
> \it{ʧaŋa}, \it{saŋa} ‘branch’; *kutu > \it{kudu},
\it{hudu} ‘headlouse’, *siku > \it{ʧiku}, \it{ʧihu} ‘elbow’,
and *baqəRu > \it{wou}, \it{ɓou} ‘new’.
Whether these differences are due to dialectal differences or
have some other basis is currently unknown.

The reflexes of *ə are particularly diverse in Bima,
with the syllable position
and qualities of adjacent vowels
and consonants all playing a role in conditioning which reflexes occur.
Sanggar usually has *ə > \it{e} while Kolo usually has *ə > \it{o}
in available data, though exceptions occur.
Before a syllable containing \it{u} Kolo
and Bima have *ə > \it{u},
and Sanggar has *ə > \it{e}.
Examples are given in \trf{03-05}.

\begin{table}[p]
	\centering\tcp{\tsc{East Sumbawa} *ə /{\gap}(C)u}{03-05}
	\st{4.25pt}\begin{tabular}{llllllll}\lsptoprule
*gloss&egg&sugarcane&three&new&PP/CC&fart&gall\\
PMP&*qat\tbr{ə}luR&*t\tbr{ə}buh&*t\tbr{ə}lu&*baq\tbr{ə}Ru&*\tbr{ə}mpu&**p\tbr{ə}su&*qap\tbr{ə}ju\\ \midrule
Sanggar&{\hr\hp{qa}}t\tbr{e}lu&{\hr}t\tbr{e}wu&{\hr}t\tbr{e}lu&{\hr}w\tbr{e}u&&{\hr\hr}p\tbr{e}su&\\
Kolo&{\hr\hp{qa}}t\tbr{o}lu&{\hr}t\tbr{o}wu&&&{\hr}\tbr{o}mpu-na&&\\
Bima&{\hr\hp{qa}}d\tbr{o}lu&{\hr}d\tbr{o}ɓu&{\hr}t\tbr{o}lu&{\hr}ɓ\tbr{o}u&{\hr}\tbr{o}mpu&{\hr\hr}p\tbr{o}ʧu&{\hr\hp{qa}}f\tbr{o}lu\\
		\lspbottomrule
	\end{tabular}
\end{table}

Bima usually (though not universally) has *ə > \it{e,
i} before syllables containing \it{i},
including when this \it{i} is not a reflex of *i.
There are two examples of Kolo *ə > \it{o} in this environment
and two of Sanggar *ə > \it{e}.
Examples of *ə before syllables containing \it{i} in Bima are
given in \trf{03-06}.

\begin{table}[p]
	\centering\tcp{\tsc{East Sumbawa} *ə /{\gap}(C)i}{03-06}
	\begin{tabular}{llllllll}\lsptoprule
*gloss	&	buy	&	low-tide	&	give	&	sneeze	&	full	&	stand	&	six	\\	
PMP	&	*b\tbr{ə}li	&	*ma-q\tbr{ə}ti	&	*b\tbr{ə}Ray	&	*b\tbr{ə}ñan	&	*p\tbr{ə}nuq	&	*k\tbr{ə}dəŋ	&	*\tbr{ə}nəm	\\	\midrule
Sanggar	&	{\hr}w\tbr{e}li	&	{\hr}m\tbr{e}ti	&	\<b\tbr{è}kè\>	&		&	\cg{(nisi)}	&	\cg{(tendi)}	&		\\	
Kolo	&	{\hr}w\tbr{o}li	&	{\hr}m\tbr{o}ti	&	{\hr}w\tbr{e}ya	&		&		&	{\hr}h\tbr{o}ru	&	{\hr}\tbr{o}nu	\\	
Bima	&	{\hr}w\tbr{e}li	&	{\hr}m\tbr{o}ti	&	{\hr}mb\tbr{e}i	&	{\hr}ɓ\tbr{e}ni	&	{\hr}ɓ\tbr{i}ni	&	{\hr}k\tbr{i}di	&	{\hr}\tbr{i}ni	\\	
		\lspbottomrule
	\end{tabular}
\end{table}

Before a syllable containing \it{a},
the reflexes are mixed,
with no clear conditioning environments,
except the presence of labial *p in *dəpa ‘fathom’
and *əpat `four' probably conditioning *ə > \it{u}.
Examples are given in \trf{03-07}.

\begin{table}[p]
	\centering\tcp{\tsc{East Sumbawa} *ə /{\gap}(C)a}{03-07}
	\st{5.5pt}\begin{tabular}{llllllll}\lsptoprule
*gloss&alone&one&hear&fathom&four&heartwood&heavy\\
PMP&*ma-\tbr{ə}sa&*\tbr{ə}sa&*d\tbr{ə}ŋəR&*d\tbr{ə}pa&*\tbr{ə}pat&*t\tbr{ə}Ras&*b\tbr{ə}Rəqat\\ \midrule
Sanggar&{\hr}m\tbr{e}sa&&{\hr}pal\tbr{i}ŋa&{\hr}nd\tbr{u}pa&&\cg{(‹ka-ekh›)}&\cg{(dodo)}\\
Kolo&{\hr}m\tbr{e}sa&{\hr}\tbr{o}sa&{\hr}pal\tbr{i}ŋa&&{\hr}\tbr{u}pa&‹kat\tbr{ò}o›&\\
Bima&{\hr}k\tbr{e}se&&{\hr}r\tbr{i}ŋa&{\hr}nd\tbr{u}pa&{\hr}\tbr{u}pa&{\hr}t\tbr{e}ra&{\hr}b\tbr{a}ra\\
		\lspbottomrule
	\end{tabular}
\end{table}

In final syllables, Sanggar usually has *ə > \it{o,} Kolo usually
has *ə > \it{u,}
and Bima usually has *ə > \it{i},
\it{a} though there are exceptions in all three languages.
Examples are given in \trf{03-08}.

\begin{table}[p]\centering\tcp{\tsc{East Sumbawa} *ə /σ\und{σ}{\#}}{03-08}
	\begin{threeparttable}\st{3pt}
	\begin{tabular}{lllll@{}ll@{}l}\lsptoprule
*gloss&stand&six&vine&inside&hear&wet&black\\
PMP&*kəd\tbr{ə}ŋ&*ən\tbr{ə}m&*waR\tbr{ə}j\su{†}&*dal\tbr{ə}m\su{‡}&*dəŋ\tbr{ə}R&*bas\tbr{ə}q&*ma-qit\tbr{ə}m\\	\midrule
Sanggar&\cg{(tendi)}&&\cg{(ika)}&‹dal\tbr{o}kh›, lal\tbr{o}k&{\hr}paliŋ\tbr{a}&\cg{(ŋgesak)}&{\hr}mut\tbr{o}k\\
Kolo&{\hr}hor\tbr{u}&{\hr}on\tbr{u}&{\hr}a\tbr{u}&‹kaḍal\tbr{u}›&{\hr}paliŋ\tbr{a}&&{\hr}kamet\tbr{e}\\
Bima&{\hr}kid\tbr{i}&{\hr}in\tbr{i}&{\hr}a\tbr{i}&\cg{(ɗei)}&{\hr}riŋ\tbr{a}&{\hr}mbeʧ\tbr{a}&{\hr}meʔ\tbr{e}\\
		\lspbottomrule
	\end{tabular}
		\begin{tablenotes}\tnsize
			\item[†]	{Reflexes here mean ‘rope’.}
			\item[‡]	{Reflexes here mean ‘deep’.}
		\end{tablenotes}
	\end{threeparttable}
\end{table}

